\chapter{From ordinary English to mathematical statements}\label{ch:english}
\begin{readergoal}Read, question, and translate statements before manipulating symbols. No mathematical prerequisite is assumed.\end{readergoal}
\section{Why begin with sentences?}
A scientist looks at an experiment and makes claims: ``this animal ran faster,'' ``these spikes followed the stimulus,'' or ``these electrodes recorded the same rhythm.'' At first, the claims are expressed in ordinary language. Mathematics begins when we decide exactly what would make a claim true, how it could be false, and what follows from it. The symbols are a shorthand for decisions already made in words. Learning a formula without learning its grammar is like memorizing a sentence in a language one cannot read.

Consider ``every recorded neuron fired.'' Does this refer to every neuron in the brain, every electrode channel, or every neuron retained after spike sorting? The statement is incomplete until we specify its domain of discourse. Consider ``some neurons were silent.'' The word \emph{some} means \emph{at least one}; it does not exclude the possibility that many neurons were silent. The word \emph{or} in standard mathematics ordinarily permits both alternatives. When a manuscript says ``the signal is oscillatory,'' we must also decide what counts as an oscillation. A label is not yet an operational test.

\section{Statements, questions, definitions, and hypotheses}
A \emph{proposition} is a declarative statement to which a truth value can be assigned, given a specified interpretation. ``Two is even'' is a proposition. ``Is two even?'' is a question, while ``Let $n$ be the number of spikes'' is an instruction introducing notation. ``For every $t$ in the recording interval, $x(t)>0$'' becomes a proposition once $x$ and the recording interval are fixed. An \emph{open sentence} such as ``$n$ is even'' has a free variable and does not yet assert anything about a particular number.

A \emph{definition} establishes how a term or symbol will be used; a \emph{hypothesis} asserts a fact that may or may not hold; a \emph{theorem} is a proposition established from previously stated assumptions by valid reasoning. In applied science we must further distinguish a mathematical theorem, which follows within a model, from an empirical assertion about the world, which requires observations.

\section{Translating ordinary English}
Let $U$ be the set of units retained after quality control. Introduce the predicate $F(u)$ for ``unit $u$ fired in a designated 100-ms window.'' Then ``every retained unit fired'' becomes
\[ \forall u\in U,\ F(u). \]
``At least one retained unit did not fire'' becomes $\exists u\in U$ such that $\neg F(u)$. These are logical negations of one another. ``No unit fired'' instead means $\forall u\in U,\neg F(u)$. Notice the difference between a negative statement about \emph{all} units and a statement that \emph{not all} units fired.

In data analysis, ``higher gamma amplitude when theta phase is near its peak'' must be turned into an explicit comparison: define phase bins, summarize amplitudes, select a test statistic, and specify a null hypothesis. This transition from sentence to mathematical object is the beginning of sound methodology.

\section{A claim is only as good as its scope}
Suppose 40 of 50 detected units increase firing rate during a task. It is legitimate to say ``80 percent of these detected units increased.'' It is not automatically legitimate to say ``80 percent of neurons in this brain area increased.'' Electrode sampling, sorting errors, and selection rules change the population to which a claim refers. We will repeatedly separate \emph{a set of observed objects}, \emph{a mathematical model}, and \emph{a physiological population}.

\begin{workedexample}[The ambiguity of ``only'']
Compare ``only neurons in CA1 show a theta rhythm'' with ``neurons in CA1 show only a theta rhythm.'' The first restricts the location of units exhibiting theta; the second restricts the rhythms observed in CA1 units. Write down what observation would contradict each statement. The two claims can disagree even on the same recording.
\end{workedexample}
\begin{exercises}
\item Translate ``not every trial contains a ripple'' into a quantified sentence and simplify its negation.
\item Explain why ``all spikes occur in bursts'' can be vacuously true for an empty spike train.
\item Identify the observational unit in each claim: neuron, trial, event, animal, and session.
\item Write one empirically falsifiable claim about coherence and one purely mathematical definition of coherence.
\end{exercises}